% Options for packages loaded elsewhere
\PassOptionsToPackage{unicode}{hyperref}
\PassOptionsToPackage{hyphens}{url}
\documentclass[
]{article}
\usepackage{xcolor}
\usepackage[margin=1in]{geometry}
\usepackage{amsmath,amssymb}
\setcounter{secnumdepth}{-\maxdimen} % remove section numbering
\usepackage{iftex}
\ifPDFTeX
  \usepackage[T1]{fontenc}
  \usepackage[utf8]{inputenc}
  \usepackage{textcomp} % provide euro and other symbols
\else % if luatex or xetex
  \usepackage{unicode-math} % this also loads fontspec
  \defaultfontfeatures{Scale=MatchLowercase}
  \defaultfontfeatures[\rmfamily]{Ligatures=TeX,Scale=1}
\fi
\usepackage{lmodern}
\ifPDFTeX\else
  % xetex/luatex font selection
\fi
% Use upquote if available, for straight quotes in verbatim environments
\IfFileExists{upquote.sty}{\usepackage{upquote}}{}
\IfFileExists{microtype.sty}{% use microtype if available
  \usepackage[]{microtype}
  \UseMicrotypeSet[protrusion]{basicmath} % disable protrusion for tt fonts
}{}
\makeatletter
\@ifundefined{KOMAClassName}{% if non-KOMA class
  \IfFileExists{parskip.sty}{%
    \usepackage{parskip}
  }{% else
    \setlength{\parindent}{0pt}
    \setlength{\parskip}{6pt plus 2pt minus 1pt}}
}{% if KOMA class
  \KOMAoptions{parskip=half}}
\makeatother
\usepackage{graphicx}
\makeatletter
\newsavebox\pandoc@box
\newcommand*\pandocbounded[1]{% scales image to fit in text height/width
  \sbox\pandoc@box{#1}%
  \Gscale@div\@tempa{\textheight}{\dimexpr\ht\pandoc@box+\dp\pandoc@box\relax}%
  \Gscale@div\@tempb{\linewidth}{\wd\pandoc@box}%
  \ifdim\@tempb\p@<\@tempa\p@\let\@tempa\@tempb\fi% select the smaller of both
  \ifdim\@tempa\p@<\p@\scalebox{\@tempa}{\usebox\pandoc@box}%
  \else\usebox{\pandoc@box}%
  \fi%
}
% Set default figure placement to htbp
\def\fps@figure{htbp}
\makeatother
\setlength{\emergencystretch}{3em} % prevent overfull lines
\providecommand{\tightlist}{%
  \setlength{\itemsep}{0pt}\setlength{\parskip}{0pt}}
\usepackage{amsmath}
\usepackage{amssymb}
\usepackage{booktabs}
\usepackage{caption}
\usepackage{float}
\usepackage{placeins}
\usepackage{booktabs}
\usepackage{longtable}
\usepackage{array}
\usepackage{multirow}
\usepackage{wrapfig}
\usepackage{float}
\usepackage{colortbl}
\usepackage{pdflscape}
\usepackage{tabularx}
\usepackage{xltabular}
\usepackage{threeparttable}
\usepackage{threeparttablex}
\usepackage[normalem]{ulem}
\usepackage{makecell}
\usepackage{xcolor}
\usepackage{bookmark}
\IfFileExists{xurl.sty}{\usepackage{xurl}}{} % add URL line breaks if available
\urlstyle{same}
% fallback for those not using the hyperref driver hyperxmp:
\makeatletter
\@ifundefined{xmpquote}{\newcommand{\xmpquote}[1]{#1}}{}
\makeatother
\hypersetup{
  pdftitle={Assignment 1},
  pdfauthor={Felix Liu},
  hidelinks,
  pdfcreator={LaTeX via pandoc}}

\title{Assignment 1}
\author{Felix Liu}
\date{}

\begin{document}
\maketitle

\section{Introduction}\label{introduction}

This dataset comes from an inter-laboratory calibration study in
analytical chemistry. Copper samples with known concentrations of 0, 2,
10, 50, and 200 \(\mu\)g/L were analyzed by seven laboratories, with
five replicate measurements obtained at each concentration in each
laboratory, for a total of 175 observations. The measured concentration
is the response variable, while the known concentration is the primary
predictor.

The goal of this analysis is to characterize the overall calibration
relationship between known and measured concentration while also
accounting for systematic differences among laboratories. I first
examine the data graphically and summarize both the overall measurement
behavior and laboratory-specific deviations. I then compare a global
ordinary least squares model with several linear mixed-effects models
that allow laboratories to differ in their intercepts, slopes, or both.
I also evaluate whether a quadratic calibration curve provides a
meaningful improvement over a linear specification. The final model is
selected based on model fit, parsimony, and the extent to which it
captures both the overall concentration-response relationship and
between-laboratory heterogeneity.

\section{Exploratory Analysis}\label{exploratory-analysis}

\begin{center}
\begin{minipage}[t]{0.62\textwidth}
\centering
\scriptsize
\captionof{table}{Summary by Known Concentration}

\begin{tabular}{rrrrrr}
\toprule
Conc. & $n$ & Mean & SD & SE & Bias\\
\midrule
0 & 35 & 1.23 & 2.92 & 0.49 & 1.23\\
2 & 35 & 3.23 & 2.36 & 0.40 & 1.23\\
10 & 35 & 10.37 & 2.26 & 0.38 & 0.37\\
50 & 35 & 48.79 & 4.87 & 0.82 & -1.21\\
200 & 35 & 192.72 & 12.21 & 2.06 & -7.28\\
\bottomrule
\end{tabular}
\end{minipage}
\hfill
\begin{minipage}[t]{0.34\textwidth}
\centering
\scriptsize
\captionof{table}{Summary by Laboratory}

\begin{tabular}{lrrr}
\toprule
Lab & $n$ & Mean Error & RMSE\\
\midrule
1 & 25 & 1.72 & 4.51\\
2 & 25 & 2.06 & 5.00\\
3 & 25 & -4.48 & 9.24\\
4 & 25 & 0.84 & 4.18\\
5 & 25 & -0.61 & 4.13\\
\addlinespace
6 & 25 & -1.35 & 7.83\\
7 & 25 & -6.11 & 10.74\\
\bottomrule
\end{tabular}
\end{minipage}
\end{center}

The concentration-level summaries show that measured concentration
generally tracks the known concentration, but the direction and
magnitude of bias change across the concentration range. At the lower
concentrations, the measurements tend to be slightly higher than the
true values in absolute terms. This is especially noticeable at 2
\(\mu\)g/L, where the absolute bias is only about 1.2 \(\mu\)g/L but
represents a relatively large error compared with the true
concentration. At 10 \(\mu\)g/L, the average measurement is very close
to the known value. At the higher concentrations, the direction of bias
reverses: measurements are slightly low at 50 \(\mu\)g/L and more
noticeably low at 200 \(\mu\)g/L, suggesting some underestimation at the
upper end of the calibration range.

The standard deviation also increases substantially at higher
concentrations, indicating greater variability in the individual
measurements as concentration increases. However, because each
concentration level contains 35 observations, the standard errors of the
mean remain considerably smaller than the corresponding standard
deviations. Thus, the mean response at each concentration is estimated
reasonably precisely even though individual measurements become more
variable at higher concentrations.

The laboratory-level summaries show additional differences in average
measurement error and RMSE across laboratories. Because each laboratory
contributes repeated measurements across all five concentration levels,
observations within a laboratory are naturally clustered and may share
systematic calibration differences. This motivates the use of a
mixed-effects model, with laboratory treated as the clustering variable.
The subsequent analysis considers whether laboratories differ mainly in
their baseline measurement level, in their response to increasing true
concentration, or in both.

\begin{center}\includegraphics[width=0.7\linewidth]{C:\Users\felix\OneDrive\Desktop\R Code\STAT 35800 Statistical Applications\Assignment 1\output\analysis_files/figure-latex/init_plot_by_lab-1} \end{center}

\begin{center}\includegraphics[width=0.85\linewidth]{C:\Users\felix\OneDrive\Desktop\R Code\STAT 35800 Statistical Applications\Assignment 1\output\analysis_files/figure-latex/spaghetti-plot-1} \end{center}

The plot shows a strong positive relationship between known and measured
concentration across all laboratories, consistent with an approximately
linear calibration relationship. At the same time, the
laboratory-specific patterns do not overlap perfectly. Some laboratories
appear to measure systematically above or below others, and the
separation among laboratories may also change as the true concentration
increases.

These patterns provide graphical evidence that laboratory-specific
effects should be considered rather than relying only on a single pooled
regression line. However, the plot alone does not determine the
appropriate random-effect structure. I therefore compare models with
different combinations of random intercepts and random slopes and
evaluate whether the additional complexity improves model fit. I also
compare linear and quadratic specifications to assess whether curvature
in the calibration relationship is substantial enough to justify
including a quadratic term.

\section{Global OLS}\label{global-ols}

As a baseline, I first fit a global ordinary least squares regression
using measured concentration as the response and true concentration as
the predictor, while ignoring laboratory-specific differences. The model
is

\[
Y_{ij} = \beta_0 + \beta_1 X_{ij} + \varepsilon_{ij},
\]

where \(Y_{ij}\) is the measured concentration, \(X_{ij}\) is the known
concentration, and \(\varepsilon_{ij}\) is the residual error. Under
this model, all laboratories share the same intercept and slope.

\begin{table}[!h]
\centering
\caption{\label{tab:OLS}Global OLS Regression Estimates}
\centering
\begin{tabular}[t]{llrrrrr}
\toprule
  & Term & Estimate & SE & 95\% CI Lower & 95\% CI Upper & $p$-value\\
\midrule
(Intercept) & Intercept & 1.071 & 0.564 & -0.043 & 2.185 & 0.059\\
conc & True concentration & 0.958 & 0.006 & 0.946 & 0.970 & 0.000\\
\bottomrule
\end{tabular}
\end{table}

The fitted calibration equation is \[
1.071 + 0.958(\text{True Concentration}).
\]

The estimated intercept of 1.071 \(\mu\)g/L represents the predicted
measured concentration when the true concentration is zero. Its 95\%
confidence interval includes zero, so the data do not provide strong
evidence that the intercept differs from zero. The estimated slope is
0.958, meaning that, on average, a 1 \(\mu\)g/L increase in true
concentration is associated with an approximately 0.958 \(\mu\)g/L
increase in measured concentration.

For a perfectly calibrated measurement process, the expected
relationship would be an intercept of 0 and a slope of 1. The estimated
slope is below 1, with a 95\% confidence interval of approximately 0.946
to 0.970. This is consistent with the pattern observed in the
descriptive summaries: measurements are relatively close to the true
values at lower concentrations but tend to increasingly underestimate
concentration toward the upper end of the calibration range. The
positive estimated intercept partially offsets this effect at low
concentrations.

\begin{table}[!h]
\centering
\caption{\label{tab:OLS summary}Global OLS Model Fit}
\centering
\begin{tabular}[t]{rrrrr}
\toprule
R.squared & Adjusted.R.squared & Residual.SD & AIC & BIC\\
\midrule
0.993 & 0.993 & 6.146 & 1136.15 & 1145.644\\
\bottomrule
\end{tabular}
\end{table}

The global OLS model explains approximately 99.3\% of the variation in
measured concentration, reflecting the very strong overall relationship
between true and measured concentration. However, this high (\(R^2\))
should not by itself be taken as evidence that the pooled OLS model is
adequate. Much of the variation in the data is naturally explained by
the large differences among the five true concentration levels.

More importantly, the global model assumes that all 175 observations are
independent and that all seven laboratories follow the same calibration
curve. The exploratory summaries and plots indicate systematic
differences among laboratories, while measurements collected within the
same laboratory are likely to be more similar than measurements from
different laboratories. The global OLS model therefore provides a useful
description of the overall mean trend, but it does not account for the
hierarchical structure of the data. This motivates the mixed-effects
models considered in the next section.

\section{Mixed-Effect Linear
Regression}\label{mixed-effect-linear-regression}

The global OLS model assumes that all observations follow the same
calibration curve and are independent. However, the measurements are
clustered within laboratories, and the exploratory analysis suggests
systematic differences among laboratories. I therefore fit linear
mixed-effects models that retain a common population-level relationship
between true and measured concentration while allowing individual
laboratories to deviate from that relationship.

In this setting, level 1 represents the individual concentration
measurements within a laboratory, while level 2 represents the seven
laboratories.

\subsection{Random-Intercept Model}\label{random-intercept-model}

I first fit a random-intercept model, which allows each laboratory to
have its own baseline calibration level while assuming that the effect
of true concentration is the same across laboratories. The model can be
written as

\[
Y_{ij} = \beta_{0j} + \beta_1 X_{ij} + \epsilon_{ij},
\]

where

\[
\beta_{0j} = \beta_0 + u_{0j}.
\]

Combining the two levels gives \[
\beta_0 + \beta_1 X_{ij} + u_{0j} + \epsilon_{ij},
\]

where \(u_{0j}\sim N(0,\tau_0^2)\) represents the laboratory-specific
deviation from the overall intercept and
\(\epsilon_{ij}\sim N(0,\sigma^2)\) represents within-laboratory
residual variation.

\begin{table}[!h]
\centering
\caption{\label{tab:random intercept}Fixed Effects for the Random-Intercept Model}
\centering
\begin{tabular}[t]{llrr}
\toprule
  & Fixed Effect & Estimate & SE\\
\midrule
(Intercept) & Intercept & 1.071 & 1.129\\
conc & True concentration & 0.958 & 0.005\\
\bottomrule
\end{tabular}
\end{table}

\begin{table}[!h]
\centering
\caption{\label{tab:random intercept}Variance Components for the Random-Intercept Model}
\centering
\begin{tabular}[t]{lr}
\toprule
Variance Component & SD\\
\midrule
Lab intercept SD & 2.673\\
Residual SD & 5.495\\
\bottomrule
\end{tabular}
\end{table}

The estimated fixed slope remains close to that of the global OLS model,
but the random-intercept model additionally allows laboratories to be
systematically shifted above or below the overall calibration curve. The
estimated standard deviation of the laboratory-specific intercepts is
2.673 \(\mu\)g/L, representing the amount of between-laboratory
variability in baseline measurement after accounting for true
concentration. The residual standard deviation is 5.495 \(\mu\)g/L and
represents remaining variation among measurements within laboratories.

\subsection{Random-Intercept and Random-Slope
Model}\label{random-intercept-and-random-slope-model}

The random-intercept model still assumes that every laboratory responds
to increasing true concentration at exactly the same rate. The
exploratory plot suggests that the calibration lines may also differ in
slope, so I next allow both the intercept and the effect of true
concentration to vary by laboratory.

The level-1 model remains

\[
Y_{ij} = \beta_{0j} + \beta_{1j}X_{ij} + \epsilon_{ij},
\]

but the level-2 equations are now

\[
\beta_{0j} = \beta_0 + u_{0j}
\]

and

\[
\beta_{1j} = \beta_1 + u_{1j}.
\]

The combined model is therefore \[
\beta_0 + \beta_1X_{ij} + u_{0j} + u_{1j}X_{ij} + \epsilon_{ij}.
\]

The random intercept \(u_{0j}\) represents how laboratory \(j\)'s
baseline measurement differs from the overall intercept, while the
random slope \(u_{1j}\) represents how that laboratory's response to
increasing true concentration differs from the overall slope. The two
random effects are allowed to be correlated.

\begin{table}[!h]
\centering
\caption{\label{tab:random intercept and slope}Fixed Effects for the Random-Intercept and Random-Slope Model}
\centering
\begin{tabular}[t]{llrr}
\toprule
  & Fixed Effect & Estimate & SE\\
\midrule
(Intercept) & Intercept & 1.071 & 0.424\\
conc & True concentration & 0.958 & 0.017\\
\bottomrule
\end{tabular}
\end{table}

\begin{table}[!h]
\centering
\caption{\label{tab:random intercept and slope}Variance Components for the Random-Intercept and Random-Slope Model}
\centering
\begin{tabular}[t]{lr}
\toprule
Random-Effect Component & Estimate\\
\midrule
Lab intercept SD & 0.439\\
Lab slope SD & 0.045\\
Intercept-slope correlation & 1.000\\
Residual SD & 4.245\\
\bottomrule
\end{tabular}
\end{table}

The estimated random-slope standard deviation is 0.045, indicating that
the relationship between true and measured concentration varies across
laboratories rather than differing only by a constant vertical shift.
The estimated correlation between laboratory-specific intercepts and
slopes is 1. This correlation describes whether laboratories with higher
baseline measurements also tend to have steeper calibration slopes.

\subsection{Comparison of Random-Effect
Structures}\label{comparison-of-random-effect-structures}

To determine whether allowing laboratory-specific slopes improves the
model beyond allowing only laboratory-specific intercepts, I compared
the two linear mixed-effects models using AIC, BIC, and the maximized
log-likelihood. Smaller AIC and BIC values indicate a preferable balance
between model fit and complexity.

\begin{table}[!h]
\centering
\caption{\label{tab:compare-random-effects}Comparison of Random-Effect Structures}
\centering
\begin{tabular}[t]{lrrrrl}
\toprule
Model & Parameters & AIC & BIC & Log-Likelihood & Singular\\
\midrule
Random intercept & 4 & 1114.52 & 1127.18 & -553.26 & FALSE\\
RI + slope, uncorrelated & 5 & 1037.12 & 1052.95 & -513.56 & FALSE\\
RI + slope, correlated & 6 & 1037.89 & 1056.88 & -512.95 & TRUE\\
\bottomrule
\end{tabular}
\end{table}

The random-intercept and random-slope model has a BIC of 1056.88,
compared with 1127.18 for the random-intercept model. The
likelihood-ratio comparison gives \(\chi^2 = 80.63\) with 2 degrees of
freedom and a p-value of \textless2e-16.

\section{Mixed-Effect Quadratic
Regression}\label{mixed-effect-quadratic-regression}

The descriptive results suggest that the amount of bias may change
across the concentration range. Measurements tend to be somewhat high at
the lower concentrations but low at the upper end, raising the
possibility that a strictly linear calibration curve may not fully
capture the mean relationship. I therefore extend the previous
mixed-effects model by adding a quadratic term for true concentration.
The model is

\[
\beta_0 + \beta_1X_{ij} + \beta_2X_{ij}^{2} + u_{0j} + u_{1j}X_{ij} + \epsilon_{ij},
\] where \(\beta_2\) represents curvature in the population-average
calibration relationship. As before, \(u_{0j}\) allows laboratory \(j\)
to have a laboratory-specific intercept and \(u_{1j}\) allows its linear
slope to differ from the population-average slope. The model continues
to assume a common residual variance across laboratories.

\begin{table}[!h]
\centering
\caption{\label{tab:quadratic}Fixed Effects for the Quadratic Mixed-Effects Model}
\centering
\begin{tabular}[t]{llrr}
\toprule
  & Fixed Effect & Estimate & SE\\
\midrule
(Intercept) & Intercept & 1.1646 & 0.4822\\
conc & True concentration & 0.9491 & 0.0294\\
I(conc\textasciicircum{}2) & True concentration squared & 0.0000 & 0.0001\\
\bottomrule
\end{tabular}
\end{table}

The quadratic coefficient determines whether the rate of change in
measured concentration changes systematically as true concentration
increases. A coefficient close to zero would indicate little curvature
and would favor retaining the simpler linear calibration relationship.

To determine whether the quadratic term provides enough improvement to
justify the additional complexity, I compare the quadratic model with
the linear random-intercept and random-slope model using AIC, BIC, and
the maximized log-likelihood. Both models are fitted by maximum
likelihood and have the same random-effect structure, so the comparison
focuses on the addition of the quadratic fixed effect.

\begin{table}[!h]
\centering
\caption{\label{tab:comparison linear quadratic}Comparison of Linear and Quadratic Mixed-Effects Models}
\centering
\begin{tabular}[t]{lrrrr}
\toprule
Model & Parameters & AIC & BIC & Log-Likelihood\\
\midrule
Linear & 5 & 1037.12 & 1052.95 & -513.56\\
Quadratic & 6 & 1039.02 & 1058.01 & -513.51\\
\bottomrule
\end{tabular}
\end{table}

Because the linear model is nested within the quadratic model, I also
perform a likelihood-ratio comparison:

The quadratic model has a BIC of 1058.01, compared with 1052.95 for the
linear model. Because smaller BIC values are preferred, the comparison
favors retaining the simpler linear specification.

The likelihood-ratio test gives \(\chi^2 = 0.1\) with degree of freedom
and \(p = 0.752\). This does not provide strong evidence that the
quadratic term improves model fit.

Taken together, the model-comparison results do not indicate that the
additional curvature provides enough improvement to outweigh the greater
complexity of the quadratic model. The linear calibration curve is
therefore more parsimonious.

\begin{center}\includegraphics[width=0.75\linewidth]{C:\Users\felix\OneDrive\Desktop\R Code\STAT 35800 Statistical Applications\Assignment 1\output\analysis_files/figure-latex/fitted-calibration-curves-1} \end{center}

\end{document}
